% !PRINTED 515
\weberpage{515}

For by the assumption concerning $\alpha, \beta$,
\[ \alpha x + \beta y = \varepsilon \]
is a unit. By multiplication with $\mu$ it follows from this that
\[ \alpha \mu x + \beta \mu y = \varepsilon \mu, \]
and since $\alpha \mu$ and $\beta \mu$ are by assumption divisible by $\beta$, it follows that $\varepsilon \mu$, and consequently also the $\mu$ associated with $\varepsilon \mu$, is divisible by $\beta$.

\begin{center}
\S.\,140.

\gesperrt{Prime functionals in the field $\Omega$.}
\end{center}

By the theorems of the preceding section we have gained the means to derive the laws of divisibility of the integral functionals in the field $\Omega$ in exactly the same way as is applied in the elements of arithmetic to the natural integers. We define as follows:

1. \gesperrt{An integral functional $\pi$ of the field $\Omega$ which is not a unit is called a prime functional if, apart from the units, it is divisible only by the functionals associated with itself. Every integral functional in $\Omega$ which has other divisors besides these is called composite.}

The concept of the prime functional is accordingly essentially dependent on the field $\Omega$. The prime functionals of one field may very well be composite in another, extended field.

If $\omega$ is an arbitrary integral functional and $\pi$ a prime functional of the field $\Omega$, only two cases are possible: either $\omega$ is relatively prime to $\pi$, or $\omega$ is divisible by $\pi$; for a common divisor of $\omega$ and $\pi$ can only be either a unit or associated with $\pi$, and in the latter case $\omega$ is divisible by $\pi$. From this follows the theorem:

2. \gesperrt{If the product $\alpha \beta$ of two integral functionals $\alpha, \beta$ in $\Omega$ is divisible by a prime functional $\pi$, then one of the two factors must be divisible by $\pi$.}
